% Generated draft: requires mathematical and source review.
\documentclass[11pt]{article}
\usepackage[margin=25mm]{geometry}
\usepackage{fontspec}
\setmainfont{Latin Modern Roman}
\usepackage{amsmath,amssymb,mathtools,unicode-math}
\setmathfont{Latin Modern Math}
\usepackage{enumitem,needspace,xurl,hyperref}
\setlength{\parindent}{0pt}
\setlength{\parskip}{5pt}
\setlength{\emergencystretch}{3em}
\newcommand{\N}{\mathbb{N}}
\newcommand{\Z}{\mathbb{Z}}
\newcommand{\Q}{\mathbb{Q}}
\newcommand{\R}{\mathbb{R}}
\newcommand{\C}{\mathbb{C}}
\newcommand{\F}{\mathbb{F}}
\newcommand{\A}{\mathbb{A}}
\newcommand{\PP}{\mathbb{P}}
\newcommand{\E}{\mathbb{E}}
\newcommand{\eps}{\varepsilon}
\newcommand{\dd}{\,\mathrm{d}}
\begin{document}
\section*{UnsolvedMath ID: 30002446}
\textbf{Source problem code:} catalogue entry
\subsection{Definitions and mathematical statement}
Let \(N\) denote the number of observed sequences and let \(p\) denote the number of model parameters. Sequential Monte Carlo denotes particle-based approximation of a sequence of probability distributions; the supplied record does not specify those distributions or the underlying statistical model.

The supplied algorithmic research question is:
\begin{quote}
Can fundamental algorithmic developments enable sequential Monte Carlo models with about \(1000\) observed sequences and \(20\) parameters?
\end{quote}
Thus the stated target has \(N\approx 1000\) and \(p=20\), with no specified tolerance for the approximation in the sequence count.

The word ``enable'' has no supplied mathematical or computational criterion. In particular, the record supplies no model class, sequence alphabet or lengths, inferential target, accuracy requirement, computational budget, or required probabilistic guarantee. Consequently, this is a research objective requiring further specification, rather than a fully quantified existence or complexity assertion. No universal quantifier over models or data sets is justified by the supplied formulation.
\subsection{Short English statement}
The question asks whether algorithmic advances can make sequential Monte Carlo practical for roughly \(1000\) observed sequences and \(20\) model parameters. The model and the meaning of practical feasibility are unspecified, so a precise success criterion cannot be recovered from the record.
\subsection{Sources}
\begin{itemize}
\item Supplied local UnsolvedMath registry record, permanent ID \(30002446\), titled ``Scalable Sequential Monte Carlo for Pathogen Sequences,'' in the Oberwolfach Reports Open Problems collection; source problem code \texttt{catalogue entry}. The record gives proposed year \(2013\) and names no proposer.
\item Supplied report citation: \emph{Design and Analysis of Infectious Disease Studies} (2014), Oberwolfach Reports, DOI \(10.4171/owr/2013/55\), pp.\ 3189--3219. Supplied source and literature-source link, labeled ``original OWR report'': \url{https://doi.org/10.4171/owr/2013/55}. The collection describes these reports as documenting mathematical workshops and their open-problem discussions.
\item Supplied literature assessment, checked on \texttt{2026-08-21}: status ``open.'' The assessment states that no general primary result establishes the requested algorithmic scaling to very large sequence and parameter counts or proves a definitive barrier. This is supplied dated curation, not independently verified; its time-sensitive status and cited claims require independent verification.
\item The registry labels the statement ``exact'' and reports that manual review confirmed the extracted statement to be self-contained, mathematically coherent, and free of LaTeX or report-layout artifacts. This supplied extraction assessment does not resolve the missing mathematical specifications identified here.
\end{itemize}
\end{document}
